\documentclass[11pt]{letter}
\usepackage[margin=1in]{geometry}
\usepackage{parskip}

\signature{Hikaru Wakaura \\ Taiki Tanimae \\ QIRI (Quantum Integrated
Research Institute Inc.)}
\address{QIRI (Quantum Integrated Research Institute Inc.)\\
1--16--3 Akasaka, Minato-ku\\ Tokyo 107--0061, Japan\\
h.wakaura@deeptell.jp,\ t.tanimae@deeptell.jp}

\begin{document}
\begin{letter}{Editorial Board \\ Applied Physics Express \\
The Japan Society of Applied Physics / IOP Publishing}

\opening{Dear Editors,}

We submit the accompanying manuscript, ``Algorithm suitability for
photonic continuous-variable quantum computers: a comparative resource
and noise-robustness study,'' for consideration as a Letter in
\emph{Applied Physics Express}. The work falls within APEX's scope of
applied photonics and quantum technologies, and it addresses a question
of direct practical relevance to the design of GKP-based photonic
processors that several groups are now constructing experimentally.

\textbf{Motivation.} Photonic continuous-variable (CV) architectures
support deterministic, room-temperature generation of large cluster
states and are a leading route to fault-tolerant quantum computing.
Their cost, however, is dominated by squeezing and by the consumption of
non-Gaussian resource states. Different quantum algorithms place very
different demands on these resources, and no unified comparison has
quantified which algorithm classes are intrinsically robust on CV
hardware and how much GKP error correction actually helps. The present
work supplies that comparison and the design rules that follow from it.

\textbf{Main findings.} We port four representative algorithms---%
control-free phase estimation, randomized product formulas, quantum
Kolmogorov--Arnold networks, and Regev factoring---to a common CV model
and compare them under photon loss (exact Fock-space Kraus channel) and
Gaussian displacement with GKP correction. Three results are central:
(i)~control-free phase estimation is the most noise-robust because it
needs no controlled unitaries and, for quadratic Hamiltonians, no
non-Gaussian gates, and the property survives for non-Gaussian
Hamiltonians; (ii)~GKP correction reduces residual error only above the
finite-squeezing floor, and whether an algorithm benefits is set by a
noise- versus systematic-limited dichotomy whose empirical envelope
constants we report with $R^{2}\!\ge\!0.89$; (iii)~against a qubit
surface-code baseline, the CV advantage is structural---Gaussian
operations are free---rather than threshold-based.

\textbf{Quantitative novelty.} The core GKP-correction model is not
assumed: it is validated against an explicit Glancy--Knill
error-correction Monte Carlo to within $0.1\%$, lifting a common
methodological concern in CV-resource papers. A two-dimensional phase
diagram and Hamiltonian ensembles confirm the predictions are not
artefacts of single operating points or single instances. We are
explicit, in a dedicated paragraph, about which results are exact, which
are validated model surrogates, and which are extrapolations.

\textbf{Experimental relevance.} The natural test bed is a GKP-encoded
photonic cluster-state processor at $10$--$15$~dB squeezing---within
reach of current laboratories---where the predicted floor crossing and
the algorithm dichotomy are directly observable; control-free
spectroscopy of a small bosonic Hamiltonian is the most accessible first
demonstration. We believe this makes the manuscript a timely and useful
contribution to APEX's photonics readership.

\textbf{Compliance.} The Letter is below the APEX length and figure
limits, the abstract is within 100 words, and the manuscript carefully
avoids the discouraged language patterns. All scripts and numerical
results that reproduce every figure and table are bundled with the
submission. The manuscript has not been published and is not under
consideration elsewhere; the authors declare no competing interests.

\closing{Sincerely,}

\end{letter}
\end{document}
